% ============================================================
% INTRODUCTION: v3, written clean 2026-07-26.
% Beats B1-B7 per PAPER_RESTRUCTURE/phase4_v3/INTRO_PLAN.md.
% Settled: no contemporary papers named here (D1); no reference to our own
% earlier unpublished draft (D2); dissociation previewed as an open
% phenomenon only (D3); no length pre-compression (D4); mechanism gets one
% line (D5).
% STYLE: no em-dashes, no AI-characteristic phrasing (see STYLE.md).
% ============================================================

% [B1: the practical question]
Community detection is a fundamental task in network
analysis~\cite{fortunato2010community, fortunato202220}. Given a graph, the goal
is to partition its vertices into groups that are densely connected internally
and sparsely connected to the rest of the graph, and the resulting partitions are
used to summarize structure, to identify functional units, and to organize
downstream computation. The task is computationally demanding, and the most
widely used methods, including modularity
optimization~\cite{newman2004finding, blondel2008fast, traag2019louvain} and
flow-based approaches~\cite{rosvall2008maps}, have running times that scale with
the number of edges rather than with the number of vertices. On the graphs to
which these methods are now routinely applied, that quantity is
large~\cite{leskovec2009community}.

Reducing the vertex set is not available as a remedy, because a community
assignment is required for every vertex; removing vertices alters the problem
rather than reducing it. Reducing the edge set is available. Graph
sparsification constructs a subgraph on the same vertex set with substantially
fewer edges while preserving properties of the
original~\cite{spielman2008graph, benczur2015randomized}. If the properties
preserved include community structure, then detection can be performed on the
sparsified graph at lower cost and without loss of quality. That premise has
been pursued for fifteen years, beginning with Satuluri et
al.~\cite{satuluri2011local}, and the resulting methods are now implemented in
standard graph libraries.

% [B2: the answer the field has, and how it was actually scoped]
Satuluri et al.~\cite{satuluri2011local} introduced local similarity
sparsification, in which each vertex retains the incident edges whose endpoints
have the highest Jaccard similarity of neighbourhoods. Evaluating four clustering
algorithms on seven networks, they reported speedups ``often in the range 10-50,
with little or no deterioration in the quality of the resulting clusters,'' and
stated that ``for at least two of the four clustering algorithms, our
sparsification consistently enables higher clustering accuracies.''

The conditions under which that result was obtained are specific. The two
algorithms for which accuracy improved are fixed-$k$, balance-constrained cut
partitioners. Accuracy is measured as agreement with external ground-truth
communities rather than by any objective computed on the graph. The reported
speedups are relative to clustering runs requiring hours on graphs of up to
$1.2 \times 10^8$ edges. The networks on which the gains are largest have average
degrees of 65, 76 and 94, and two of them are derived from high-throughput assays
known to contain false-positive edges. The same paper reports a synthetic sweep
in which the method ``actually outperforms the original clustering starting from
degree 50.'' To our knowledge, no subsequent work examines that threshold.

% [B3: the scope drift, stated without accusation]
Subsequent work extended the recommendation beyond these conditions.
Sparsify-then-detect has since been applied to modularity optimizers as well as
to cut partitioners, to sparse graphs as well as to dense ones, and evaluated
against graph-internal objectives as well as against external references. The
degree threshold reported in the original sweep does not appear in this later
literature, and the conditions attached to the original result are not generally
restated with it.

% [B4: the lost control and what its absence costs (THE HINGE)]
A measurement requirement is also frequently omitted, and the omission has a
direct bearing on what can be concluded. A sparsifier intended for community
detection removes edges lying between communities in preference to edges lying
within them. Quality measures for a partition, including modularity and
conductance, penalize exactly the between-community edges that such a sparsifier
removes. A partition evaluated on the sparsified graph is therefore evaluated
against a graph from which its penalty term has been partially deleted, and the
resulting quantity reflects the sparsifier's selectivity rather than the extent
to which community structure was preserved. The comparison that answers the
intended question evaluates the partition obtained from the sparsified graph and
the partition obtained from the original graph on the same graph, namely the
original one.

This requirement is stated in the 2011 paper, whose authors write that they
``cannot simply measure the conductance using the very same sparsified graph,
since that would tell us nothing about how well the sparsified graph retained the
cluster structure in the original graph,'' and who report every quality figure on
the original graph accordingly. It is absent from much of the work that followed,
and from the released evaluation code of a widely used sparsification benchmark.

The effect of the omission is large. Across 63 controlled configurations spanning
four detection algorithms, two sparsifier families and seven networks, evaluating
on the sparsified graph rather than on the original reverses the sign of the
conclusion in 45 of them. In the most extreme configuration, a pipeline that
reduces modularity by 0.40 under the correct evaluation appears to increase it by
0.39 under the incorrect one. Under such an evaluation a boundary cannot be
located, since improvement is reported throughout the parameter range.

% [B5: what we did]
We therefore rebuilt the evaluation before re-asking the question. Partitions are
always scored on the original graph. Comparisons against a stochastic baseline
are given the same wall-clock budget as the pipeline they are compared with, and
are also checked against a fixed-restart baseline, because a runtime-matched
control is weak precisely when the pipeline is fast. Because sparsification
fragments partitions, and because best-match and information-theoretic recovery
scores rise mechanically with the number of clusters, every recovery comparison
is matched for granularity against a partition of the original graph at the same
cluster count, and is referenced to a chance floor computed from size-matched
random partitions. Because a heavy-tailed degree sequence alone can produce
several of the effects attributed to community structure, the central mechanistic
claims are checked against degree-preserving randomizations of the same graphs.

With that instrument we ran 23 controlled experiments over both sparsifier
families that have ever claimed quality gains, degree-based and similarity-based,
the latter being the 2011 method itself; four detection algorithms drawn from two
families; seventeen real networks of up to 25 million edges; and a synthetic
sweep of the one variable the original authors identified as decisive. The design
locates a boundary rather than confirming or refuting a claim, and it is reported
as such. Predictions and stopping rules were fixed before each experiment ran,
and several of our own hypotheses did not survive them.

% [B6: the boundary]
Our results locate a boundary, and the transition across it is abrupt.

On one side of it sparsification does not improve detection. For modularity-family
detectors with free granularity, on graphs with average degree between 5 and 33,
no sparsifier we tested improves modularity measured on the original graph beyond
the variation of a small number of random restarts, at any retention level that
preserves quality, under any of the four algorithms examined. Nor is there a
speed benefit. Measured end to end, including the cost of the sparsifier itself,
the pipeline requires between $0.87\times$ and $1.35\times$ the time of clustering
the original graph directly, and at aggressive retention it is slower, because a
fragmented graph presents the optimizer with more work rather than less. Every apparent
gain we found in this regime, in modularity, in ground-truth recovery, and in the
one case where a sparsifier's edge-selection signal is demonstrably
structure-aware rather than degree-driven, dissolves when granularity is matched,
when chance is subtracted, or when the baseline is given an equal compute budget.

On the other side of the boundary sparsification does improve detection, under
the conditions the original authors described. On synthetic graphs with planted communities, holding the number of
communities fixed by construction so that no granularity effect is possible,
similarity-based sparsification raises both the objective measured on the
original graph and chance-corrected agreement with the planted labels once the
average degree passes approximately 50. The effect is absent at average degree 11
and 25, present in two thirds of configurations at 49, and present in all of them
at 102 and 221. It survives the worst case over ten independent runs of the
partitioner, and it is not reproduced by degree-based or uniform sparsification
at identical retention, so it is a property of the similarity signal rather than
of the edge budget. Injecting spurious edges makes the gain grow monotonically,
which is the signature the denoising explanation predicts.

What decides the outcome is not which sparsifier is used. It is the average
degree of the graph, the family of the detection algorithm, since fixed-$k$
balanced partitioners gain where free-granularity modularity optimizers do not,
and whether quality is measured against an external reference or against an
objective defined on the graph whose edges were deleted. The mechanism is
consistent with this. The quantity governing a degree-based sparsifier's effect
on modularity is the separation between the scores it assigns to within-community
and between-community edges, and we show by direct intervention that this
quantity is controlled by where high-degree-product edges sit relative to
community boundaries, at fixed degree sequence, fixed partition and fixed
modularity.

Two qualifications attach to this result. Sparsification repairs a weaker
detector rather than surpassing a stronger one: a resolution-tuned modularity
optimizer on the untouched graph outperforms every sparsified pipeline we
measured. And with an exact similarity computation the quality gain is not a
speed gain, because the sparsifier costs more than the detection it accelerates.

% [B7: contributions]
This paper contributes the following.

\begin{itemize}
  \item \textbf{The boundary.} A degree threshold near average degree 50 and an
  algorithm-family split, established under controls that exclude granularity by
  construction, together with the negative territory delimited across two
  sparsifier families, four detection algorithms and seventeen networks.

  \item \textbf{An evaluation protocol, and the measured cost of omitting it.}
  Each control is motivated by a specific failure it catches, and the consequence
  of the central omission is quantified at 45 sign reversals in 63
  configurations. We also introduce a permuted-weight control for
  similarity-weighted pipelines, which shows that the benefit attributed to
  similarity weighting is reproduced, and in one case exceeded, when the weights
  are randomly permuted across edges.

  \item \textbf{A mechanism for the boundary.} A decomposition of the
  score-separation quantity into a sorting term and a granularity term across
  seventeen networks, and a direct causal test in which that quantity is steered
  through more than a full sign change while degree sequence, partition,
  intra-edge fraction and modularity are all held exactly fixed.

  \item \textbf{An open phenomenon.} In several controlled settings the
  partitions that best recover reference communities are not the partitions that
  score best on the objective, and the two criteria move in opposite directions
  within the same configuration. We report this as an observation rather than a
  result, and argue that it is the most consequential open question the boundary
  exposes.
\end{itemize}
